% 原书分章PDF第22—24页：第30.8节全部5段结论，TMLE因果解释及两个完整括号说明。
\section{结论}

% S30.8.P01
本文试图具体呈现所说的真理困境。真理是“人是万物的尺度”的一部分，在认识论中具有关键地位，因为包括显然更为时髦的因果性在内，较不根本的概念都预设了它。我们突出讨论了真理的若干方面，尤其包括它在险恶的大世界中的脆弱性，以及它在统计学发展全程中与这门学科的问题关系：从十九世纪大众文化中的反弹，到如今几乎所有科学都发生概率转向之际的不信任与怀疑。我们还描绘了模型概念受到的侵蚀，以及当今统计实践的折衷主义——尽管教科书错误地让人以为这是一个统一的领域——还有同样重要的、正在兴起并于数据科学中日渐流行的反统计学立场。

% S30.8.P02
前文表明，TMLE 的某些方面确实以系统方式回应了这些问题。此外，我们概述了这种反统计学立场如何可能转变为反认识论立场；后者似乎比单纯的“理论终结”更为广泛。我们看到，认识论与研究方法论的调和在许多方面都存在问题，而且不大可能仅从科学哲学中找到解决方案。从古典怀疑论、不断变化的真理定义，到逻辑实证主义对发现语境与证成语境之区分的执着、外在主义历史研究、文化相对主义和社会建构主义，所有这些立场都挑战了科学的理性与客观性，或至少强化了真理概念的衰落，或宣告它是多余的。研究方法论有着不同的、主要偏重实用的优先事项，这并不令人意外。随后，我们说明，迄今试图通过弥合统计学与科学哲学之间的鸿沟，来抵制这种发展的努力，仍存在问题。我们把这一点推广为对追求真理中一幅令人遗憾的三联画的勾勒：层次I、II与III之间，往往存在不必要的界分。由于如今围绕真理概念的一组问题，常以间接方式显现，并持续出现在科学哲学的当代主题中，我们从历史与哲学视角聚焦于因果性概念。因果概念及其术语，仍在科学语言中发挥关键作用。它在层次I、II与III中都占据主导地位，但迅速增长的因果性文献，有力地显示了前面描绘的那种不幸的分隔。因果性在哲学中有着悠久传统，经历过许多概念变化，却始终依赖语境，与认识论或形而上学立场紧密相联。因此，它一次又一次地被重新发明，受到相关思想家的形而上学与认识论观点的影响，并与之紧密相连。由于哲学中的理论或观点很少被证伪和取代，许多相互冲突的思想延续下来，演变为当今关于因果性的多元观点，似乎不容许任何精细的统一化。

% S30.8.P03
这显然给认识论带来挑战；许多哲学问题在这里甚至尚未列出，更谈不上解决。事实是，许多古老的形而上学与认识论内涵仍在发挥作用。不过，正如 Starmans（2016a，2017a）所指出的，任何解决因果性周围的哲学与实践问题的尝试，都应该考虑其当前状态；而这一状态在很大程度上由思想史中的若干重要转变决定。这些转变包括卡尔·皮尔逊与罗纳德·费希尔的概率革命——没有它，因果性概念在研究方法论中根本不会有任何地位；还包括尼尔斯·玻尔对量子力学的哥本哈根诠释，它对因果性思考产生了决定性影响；此外，还有克劳德·香农的通信理论。他关于信息是在语境与流动中的数据的思想，以及由此带来的温和革命，超越了唯物主义，使信息不再是一种附带现象或派生概念；考虑到本文讨论的主题，这一点似乎尤其值得关注。与 Pearl 这样的有力倡导者相对，我们在整个历史中也看到一种消除这个术语的倾向，代表人物包括 Pearson、Mach、Russell、Churchland、Anderson 和 Havely。不过，Google 也采取过不同立场，例如在“自动统计学家”（automated statistician）项目中（Starmans，2015b）：恰恰是在由庞大而复杂的模型自动分析成千上万个变量时，人们只有在理解机制、论证——包括因果结构与原因—结果关系——的前提下，才能对统计结论以及基于这些结论的决策保有信心。由此再次表明，解释、因果性与（自然）规律这些经典主题，在这里不应被分开处理。此外还须注意，鉴于针对统计学的反弹、它与伦理学之间令人棘手的关系，以及对道德维度和“人是万物的尺度”的诉求，数据科学的“阴暗面”如今也迅速受到关注，例如前面提到的那本新著（O'Neil，2016）。

% S30.8.P04
正如 van der Laan 和 Rose（2011）所指出的，TMLE 以不同方式处理因果性问题。因果性服从于真理与推断，纯粹统计目标参数与因果效应之间有严格区分。要在 Neyman–Rubin 因果网络之下赋予统计参数以因果解释，需要若干假设，例如正值性（即在各分层内，每个实体接受任一种暴露的概率都应非零）、无混杂因素假设，以及所谓的稳定单元处理值假设（stable unit treatment value assumption；即一个实体的暴露状态不应影响其他实体的潜在结局，而且对于所有处于某一暴露水平的个体，该暴露水平都应相同）。其重点首先在于因果估计，明确表述所关心的因果问题，并估计定义良好的统计参数；以研究问题为引导，使用对函数形式作出尽可能少假设的灵活算法，从而能够恰当地调整混杂。

% S30.8.P05
许多微妙之处在这里仅略有涉及，或不得不完全搁置。例如，承认科学事业的社会方面与外部方面，并不意味着滑向相对主义或社会建构主义；即使在实验物理学中也是如此，物理学家 Peter Galison 已在许多场合指出这一点（Galison，1987；Starmans，2017c）。同样的道理也适用于哲学与统计学之间的问题关系、数据分析中整合且统一的进路的缺乏，以及反驳前述科学“新观念”的必要性。总的说来，可以认为，上面围绕真理、因果性、数据科学与知识三联画及其问题关系所提出的议题，不仅是认识论与统计学的重大挑战，也是迈向真正的数据科学哲学的第一步。
